\documentclass{article}

% Tables live in ../tables/ (one file each, generated by analysis/paper_tables.R and the R scripts in analysis/).
% Bibliography: put sample.bib next to this file.

\usepackage[english]{babel}
\usepackage[letterpaper,top=2cm,bottom=2cm,left=3cm,right=3cm,marginparwidth=1.75cm]{geometry}
\usepackage{indentfirst}
\usepackage{float}
\usepackage{amsmath}
\usepackage{amssymb}
\usepackage{booktabs}
\usepackage{graphicx}
\usepackage{setspace}
\usepackage[authoryear, round]{natbib}
\usepackage[colorlinks=true, citecolor=blue, linkcolor=blue, urlcolor=blue]{hyperref}

\newcommand{\tbl}[1]{\input{../tables/#1}}

\title{Do Hitters Have Complementary Skill Sets That Enhance Each Other's Output?}
\author{Jai Correa}

\begin{document}
\doublespacing
\maketitle

\begin{abstract}
This paper uses hitter trait statistics to determine whether or not complementarities exist between pairs of hitters for the 2025 and 2026 seasons that boost run expectancy. This is done by establishing a two-way fixed effects model on the player and season, using standardized traits. To establish complementarities, we develop 12 hypotheses for both focal and conditional run expectancy across 2025 and 2026, the out-of-sample season. The findings from 2025 were unable to be replicated in 2026. However, we gained insight into focal delta run expectancy mechanically absorbing much of the plate-appearance production and used Barrel/PA as a proxy. Furthermore, we find that even in the case of the most reliable traits, the coefficients are sporadic since only a few pairs drive a large number of the plate appearances used in the dataset, resulting in a future study likely needing a larger sample size to gain more well-defined pairings.
\end{abstract}
\newpage
\section{Introduction}

Baseball lineups are best described as an assembly line. The end product is a culmination of runs with nine seemingly independent inputs. The manager is tasked with composing a well-oiled machine with the inputs at his disposal. In a reversal from old convention of fast players at the top of the lineup with home-run hitters occupying the middle, more recent lineups are constructed by having the best hitters occupying the top four spots of the lineup, where the best is often in the second spot (\cite{tango2007book}; \cite{bukiet1997markov}).

The Book's (\cite{tango2007book}) lineup construction comes from using a player's season-long statistics -- e.g., batting average, on-base percentage -- to simulate what a lineup would produce with different player types moved around within it. However, the season-long statistics lock in a hitter's skill set regardless of who is hitting behind or in front of them. What if certain hitters performed better or worse depending on who they were placed next to? Furthermore, it could mean heterogeneity of particular skills that boost the run-scoring ability of each hitter depending on the context of hitters before and after -- complementarity.

One of the few studies that covered hitters being complements or substitutes showed that a hitter's batting average increased by the other hitter's productivity but decreased when his team's pitchers produced more (\cite{gould2007interactions}). Rather than investigating how a player performs when dropped into a particular environment, we're asking a different question. It's whether particular hitters complement each other, rather than evaluating the gain in one hitter from the environment they were placed in. Furthermore, it's hard to justify every hitter affecting a particular hitter. For example, it seems highly unlikely that the eighth hitter in a lineup impacts the production of the fourth hitter, and vice versa.

This paper examines which hitter skills influence the production of the hitter either directly preceding or following them.  The idea is that a preceding hitter's particular skill -- say, how often they chase outside of the strike zone -- and a subsequent hitter's secondary skill -- e.g., how often they make contact -- work together to boost the run expectancy of either the preceding or the subsequent hitter.

Using basketball could be an interesting vessel to investigate hitters' complementarities. In basketball, players are complementary on the floor since there is one basketball and each presents an option to score. That means where one player shoots well from opens up space for another player to score from their strong point (\cite{kuehn2016complementary}). While we would have to account for baseball's rigid lineups instead of basketball's fluid ones, which offer greater variety in groupings, the focus on the cohesiveness of player traits is worth investigating.

Using this framework, we first construct hypotheses based on the compatibility of hitters' traits in boosting run production for both hitters during the 2025 season, broken into two arguments. First, if both hitters' traits move in the same positive or negative direction, there is a direct relationship between sign and run expectancy.  Using bat-tracking traits, we also assign each trait both a positive and a negative connotation. For example, a high swing path tilt (i.e., a steeper uppercut) would imply more frequent power, but also more frequent swing-and-miss. Second, we use contrasting style traits -- absolute contact gap and nitro zone distance -- with the expectation that a greater contrast makes it harder on the pitcher to adjust, leading to greater success for run expectancy. We end up finding a few significant interactions in 2025, some opposite to our hypothesis, though none of them are validated in 2026.

\section{Data}

The data collected is from Statcast, Major League Baseball's publicly available data. The first set of analyses is done on the 2025 season before using 2026 as robustness checks. We also only include players that have seen at least 50 plate appearances during each season so we establish a solid baseline for each player. We also use only regular-season data, excluding Spring Training and the Postseason. Furthermore, plate appearances that end an inning in the middle of the at-bat -- which include a caught stealing with two outs already in the inning or a wild pitch that scores the winning run -- have been removed from the dataset so that there are no self-joined pairs. We ensure this because if a caught stealing ends the inning, the hitter gets to hit again to lead off the next inning.

The data formed is a mixture of pitch-level and plate-appearance (PA) level statistics. PA-level includes traditionally known statistics such as batting average, on-base percentage, and slugging percentage. In this study, there are also two PA-level statistics, the first being Barrels/PA. Barrels are MLB's designation for a well-hit batted ball, based on an exit velocity of at least 98 MPH and a launch angle off the bat between 26 and 30 degrees. For each additional one MPH off the bat, the range grows one degree in either direction until 116 MPH is reached (results in a launch angle between eight and 50 degrees). A Barrel often results in a minimum batting average of .500 and a minimum slugging percentage of 1.500 (\cite{mlb_barrel}). For our discussions, a Barrel results in damage so it's a good trait in evaluating a hitter's power without the luck-induced results. Included in the data is also a cartesian location of these barrels to help identify a 'nitro zone' of that particular hitter, which will be identified in the following section. The second is Sweet Spot\%. This is MLB's designation of the optimal launch angle for batted balls, falling between eight and 32 degrees (\cite{mlb_sweetspot}). This excludes exit velocity, as soft line drives still end up as hits a large number of the time.

Pitch-level statistics include examples such as how often a hitter sees a pitch is in the zone, chases a pitch out of the zone, and swings and misses. These give us a better indication of what decisions a hitter is making over the course of the season.

In the pitch-level bucket, we also used Statcast's bat-tracking metrics in here to show how a hitter swings the bat. This includes a few metrics worth describing, as shown below:
\begin{itemize}
    \item Attack angle median, which is angle of the bat as it makes contact with the ball. A steeper attack angle often speaks to power hitters with a steep upper cut through the ball (\cite{mlb_attackangle}).
    \item Swing path median tilt, which is the bat's angle from the horizontal line at the shoulders from the pitcher's viewpoint. This is a similar idea to the attack angle in showing power swings, except this shows body tilt rather than just bat tilt (\cite{mlb_swingpathtilt}).
    \item Attack direction median, which is the angle of the bat approaching the ball horizontally. The direction is positive when going to the hitter's opposite field, and negative when going to the pull side (\cite{mlb_attackdirection}).
\end{itemize}

\subsection{Forming the Pairings}

We call the first hitter the conditional and the second hitter the focal. The conditional and focal hitters are paired based on their respective game at-bat number and team played for. Team played for is used because Statcast records the at-bat number constantly ticking through the game, rather than splitting amongst teams. This means without the team kept in mind, we would get players of two different teams often getting paired together (e.g. last of top half of inning and first of bottom half of inning). As mentioned earlier, we also remove truncated plate appearances to prevent a hitter from being joined on himself.

Once we form the pairs, we can record both the conditional and focal hitter's pairing production, as well as their season-long averages. We record these in each row of the data set, including the number of at-bats for that pairing. We also split the data by the hitter's handedness. While switch-hitters are not in the majority, they do have a meaningful impact on this study, since we are studying how the skill sets pair together. For example, Cal Raleigh of the Seattle Mariners is a switch hitter who has hit behind right-handed Randy Arozarena. Instead of grouping all of Raleigh's at-bats behind Arozarena, this dataset splits them into one row for Raleigh hitting left-handed behind Arozarena and another for Raleigh hitting right-handed in that scenario.

After cleaning our data as told above, we achieve the following distribution of pairs.

\tbl{pair_pa_distribution}

In each season, over a third of pairings had fewer than 30 plate appearances together, and half had fewer than 50. To keep the larger sample size at play, we include these pairings at play knowing that the regression will likely have a greater weight on the greater PA pairings anyway.

However, our pairing setup has a limitation: the leadoff hitter is not a focal hitter, so his plate appearance is dropped. Furthermore, a player who ends the game is a focal hitter, so there is no offsetting measure. We can overlook this since we are only losing one plate appearance, and a team averages 38 plate appearances (around 3\%).

\subsection{Understanding the 'Nitro Zone'}

The concept of the Nitro Zone is developed by the author, where the barrels of a player come in a highly concentrated area of the zone, rather than being sporadic (\cite{correa2021dalbec}). The idea is to highlight the hitter's strengths in the strike zone. Rather than constructing a heatmap and defining a hitter's 'nitro zone', we used the median Cartesian point of a hitter's barrels for simplicity. In the dataset, 'nitro x' and 'nitro z' refer to the horizontal distance relative to the center of the plate and the vertical distance off the ground, respectively.

\subsection{Delta Run Expectancy}

The delta of the run expectancy is our dependent variable of choice. It is an average of runs scored for each of the 24 base-out states (RE24) -- eight possible combinations of baserunners for each out. Statcast defines the variable as a change in run expectancy per pitch, which is another wrinkle in the matter. This means that we will sum the run values over each PA, to include the changes to the result with respect to the count. For example, a two-strike home run hit would net a larger run expectancy jump than one hit on a 3-0 count.

For reference throughout the remainder of this paper, the mean delta run expectancy per full-PA is 0.008 in 2025 (0.009 in 2026), so we should expect small coefficients in each of our regressions.

\section{Methodology}

\subsection{Model}

Using the delta of run expectancy as our dependent variable (of either the conditional or focal hitter), we seek to find which particular traits of both hitters in the pair by running a two-way fixed effects model. Given the nature of our data, the regression is only identified on the defined pairings with hitters seeing at least 50 plate appearances in 2025 (and in 2026). Furthermore, the coefficients are only causal if there aren't unobservables also affecting performance.

$$y_{fcs} = \beta_1 \, X^{cond}_{c} + \beta_2 \, X^{focal}_{f} + \beta_3 \left( X^{cond}_{c} \times X^{focal}_{f} \right) + \alpha_f + \gamma_c + \delta_s + \varepsilon_{fcs}$$

Our independent variables $X^{cond}_{c}$ and $X^{focal}_{f}$ are standardized before being implemented into the regression for ease of interpretation. We want to say that a high first-trait hitter and a high second-trait hitter affect run expectancy, rather than being bogged down by minute percentages that we were originally presented with.

$\beta_3$ represents our complementarity coefficient between the two hitter's traits.  $\alpha_f$ and $\gamma_c$ represent the fixed effects for the hitter, while $\delta_s$ is the fixed effect for the season.

Standard errors are clustered on both the conditional hitter and the focal hitter. We do this since each focal hitter has many different conditional ones, and vice-versa, so there are overlapping hitters across the many pairs. To therefore combat the potential correlated errors from pairs sharing hitters, we cluster on both players.

\subsection{Expectations}

There are two groups for our hypotheses, and the explanation for the results depend on whether we are highlighting the conditional or the focal hitter.

Our first group is keeping the same direction across both hitters' traits. Below, for shorter notation, we take the first trait listed as the conditional hitter's, and the second as the focal hitter. Working exclusively in the positive direction, we use regressions (1) On-Base Percentage (OBP) x Slugging Percentage (SLG), (3) Sweet Spot\% x Attack Angle, (4) Line-drive\% x Swing Path Tilt, (5) Chase\% x Attack Direction, (8) Attack Angle x Chase\%, (10) Contact\% x Barrel/PA, (11) Chase\% x Swing\%, and (12) Pitch/PA x Swing Path Tilt. In the negative direction, (6) Swing Path Tilt x Swinging Strike\% and (9) Swing Path Tilt x Attack Angle.

Most of the positive expectations in that group are fairly self-explanatory. (1) is simple -- a player gets on base, a second hitter drives him in with an extra-base hit. In the reverse direction, we still expect positive run expectancy because a player who can get on base likely walks and, more importantly, stays in the strike zone more often, and pitchers will not want to fall behind to that hitter when a power hitter is on deck. With pitchers likely throwing more pitches in the heart of the plate to avoid this, we expect the conditional hitter to also boost their run expectancy in those particular pairings. (3), (4), (5), (8), and (10) follow a similar notion, where a conditional hitter that either hits the ball well or makes frequent contact in front of a power hitter results in more runs in both directions. As a reminder, we still expect a positive run expectancy in both directions for (5), (8), and (11) despite having our expected negative interaction sign. Both interactions have Chase\%, and a positive hitter trait is associated with a lower Chase\% when the focal trait is high, resulting in a negative coefficient.

The set of negative directions shows how we change the connotation of bat-tracking metrics. While an increased Swing Path Tilt and Attack Angle can result in hitting more power, it also increases the risk of swinging and missing. (6) and (9) can then be interpreted as two high-whiff hitters that would decrease run expectancy in both directions.

The second group is based on contrasting styles, with regressions (2) Absolute Contact Gap x SLG and (7) Nitro Dist x Attack Direction. For our conditional hitter traits, both are differences between the two hitters in the pair, showing how the conditional hitter's difference influences run expectancy.

The contact gap in (2), regardless of direction, represents two different hitting styles between the players. If you have a player that makes less contact than the following one that has a high SLG, it would likely help the former hitter see better pitches to hit that are easier to make contact on, boosting the former hitter's run expectancy. Working in that same loop also then boosts the focal hitter's run expectancy, as a previous hitter seeing better pitches gives the power hitter a likely at-bat with a runner on base. Conversely, if the previous hitter makes more contact, the easier pitches for the high-SLG following hitter follow the same mantra as our group-one positive section, boosting both hitters' run expectancy.

Nitro Dist in (7) is used to discern the gap between hitters' strengths within the zone. Intuitively, most players hit more authoritatively to the pull side, so as Nitro Dist increases, we are more likely dealing with players hitting on opposite sides of the plate. Let's take a moment to imagine how pitchers throw the ball. They establish routines and channels of where they are throwing within the zone, and adjusting plate-appearance to plate-appearance is difficult. Hypothetically, let's say a right-handed hitter goes first with a left-handed hitter after. If he keeps away to the righty, then he's likely to throw in that same spot to the lefty after, which is now inside to him. In that case, each player is pulling the ball, which has a negative attack direction, giving us our expected negative sign for the focal run expectancy. In the opposite direction, a larger Nitro Dist plays better for a conditional hitter when the focal hitter is a pull hitter, because the conditional hitter is more likely to damage when he's pulling the ball -- a greater distance from the focal hitter's pulling strength. As you see in the table below, there are a larger number of plate appearances with opposite-handed hitters than same-sided ones in our data set, which gives us our negative signs.

\tbl{pair_side_composition}

\section{Results}

\subsection{First Series Of Regressions}

Table 5 shows the first series of regressions running on the focal hitter's delta run expectancy for 2025. In it, regressions (9) and (12) return statistically significant, except both interactions have signs opposite to our prediction. In addition to our prediction for (9), the two negative connotations for our hitters' traits also suggest how pitchers might adjust from one hitter to the next. You would think that, in the focal direction, a pitcher who has already faced a similar swinging hitter would adjust to get the focal hitter out, which makes the positive coefficient all the more surprising. (12) is an extension of (1), as a high-Pitch/PA hitter is more likely to walk, ahead of a high-Swing Path Tilt hitter who is prone to hitting for extra bases to drive him in. In the focal direction, it means that a hitter with a high-Swing Path Tilt was less successful when he was likely hitting with a runner on base. In both cases, the magnitude of each was roughly a third of the mean delta run expectancy.

Table 6 features the same regressions for 2025, except on the conditional run expectancy. Of the three statistically significant interactions, only (2) obtained the same sign as the one we predicted. The only possible explanation for (6) is that, given a low-whiff hitter behind him, a conditional hitter who also has a low Swing Path Tilt boosts his own run expectancy. It is certainly plausible, as both traits are viable proxies for high-contact hitters, which would boost run expectancy. Though, given the less imposing threat of a low-whiff hitter, a high-Swing Path Tilt conditional hitter should also benefit from being pitched around more carefully and thus increase his own run expectancy. (12) again brings us back to our Swing Path Tilt discussion of (6). As expected, because a high-Swing Path Tilt focal hitter is more reflective of a power hitter, a conditional low-Pitch/PA hitter seems to have his run expectancy boosted. A low-Pitch/PA hitter is expected to swing early and often, so if the pitcher is more concerned with the hitter on deck, those early-count pitches could be very hittable, resulting in positive run expectancy. By contrast, for the low-Swing Path Tilt focal hitter, our results show that a high-Pitch/PA conditional hitter's run expectancy is boosted. This makes sense working in the focal run expectancy direction, but in the opposite direction makes less intuitive sense. A more contact-oriented focal swing would seemingly not offer the same protection from the pitcher to the conditional as a loftier swing would. We'll table the discussion of (6) and (12) for now and leave it for our validation check on 2026.

\subsection{Robustness Check}

To verify the findings we had from the 2025 season, we use 2026. As shown in Table 7, the results of (9) and (12) could not be replicated. Given that sample sizes are roughly the same each season, the results of (9) and (12) appear to be a 2025 anomaly, making our focal hitter's run expectancy findings seem inconclusive.

Table 8 tells us that we can't replicate the findings of regressions (2), (6), and (12) in 2025, while also gaining new results in (3) and (5). Both of 2026's discoveries flipped signs from 2025. Similarly to our last robustness check, our second round of regressions regarding the conditional hitter's run expectancy are also likely inconclusive.

\subsection{A Potential Proxy For Testing The Focal Hitter}

Returning to Table 5, a clear story emerges. See our first regression: a conditional hitter's OBP with a focal hitter's SLG. That should undoubtedly result in positive run expectancy, except that its interaction term is actually negative and not statistically significant. This takes us back to the construction of our dependent variable, which is the second hitter's run expectancy minus the previous hitter's. The difference, which happens predominantly at the plate-appearance level, therefore takes each hitter's results and finds the net run expectancy between the two. Ergo, focal delta run expectancy seems to mechanically absorb production results, which is why (9) and (12) were among the hits for the 2025 season since they don't include them.

To find a proxy for our run expectancy variable, we can look at two variables: our aforementioned Barrel/PA and weighted on-base average (wOBA). wOBA is calculated using a series of linear weights for plate-appearance outcomes (e.g., single, double, walk, etc.) to then produce a number between 0 and 1, therefore reading similarly to how you would read on-base percentage. To identify a better proxy for focal run expectancy, we present the two OLS regressions below.

\tbl{proxy_regressions}

wOBA clearly explains more of our run expectancy variable than Barrel/PA does; however, there are reservations. wOBA's linear weights for each type of event are calculated by averaging the expected runs scored across each of the 24 base-out states (\cite{tango2007book}). This means that the only difference between the delta run expectancy and wOBA is that the former relies specifically on the game state, while wOBA is an average across states. Therefore, we'd likely see the same results as the first series of regressions. We would then seem to be better off using Barrel/PA, even though it is an imperfect measure, because it's not mechanically tied to run expectancy.

Now using focal Barrel/PA as our dependent variable, we see the results in Table 9. There were three statistically significant results this time around, each of which obtain the same sign we predicted they'd have. (3) and (4) fit the narrative of good hitters hitting in front of power hitters boosts the latter's Barrel/PA. Similarly, a low-chase\% conditional hitter helps a high-swing\% focal hitter boosts his Barrel/PA as well. Both are instances of gains made when hitting with a runner likely on base.

In Table 10, we again cannot replicate the 2025 findings, as shown in (3), (4), and (11). In both (3) and (4), the sign flipped from positive to negative in 2026, contrary to expectations, while (11) remained negative. However, we found a new result in (8) for 2026, reflecting the expected negative sign. This would be worth investigating across other seasons as validation.

\section{ Discussion and Limitations}

Across both 2026 validations, we were unable to replicate the findings from 2025, regardless of the particular hitter's run expectancy. Part of this likely stems from the design limitations and the traits we observed.

For one, managers choose lineups, so pairings with a larger sample size are starting players who frequently occupy a lineup. While most rows of the data reflected pairings with few plate appearances together, the regressions weight the more frequent pairings in the analysis. With a substantially smaller sample size, variation is limited, especially in one-season samples like the one we used here. Drawing closer to the start of Statcast's database with batted ball metrics (2015) would give us a larger number of pairs with a solid sample size. One caveat is that we would then lose the bat-tracking metrics, since they were only tracked from the middle of the 2023 season onward.

Another limitation is the noisiness of the traits that we have:

\tbl{trait_reliability}

Using the Spearman-Brown correction, we find that the reliability of bat-tracking and pitch-level metrics is quite stable even at lower sample sizes. However, batted-ball metrics and plate-appearance-level statistics are less reliable and require a larger sample size. This helps explain results like the drastic (+0.0008 in 2025 to -0.0016 in 2026) sign change in (4) from Tables 5 to 7, as LD\% is the least reliable trait we used in this study.

\section{Conclusion}

This paper aimed to determine which skills improved the production of a pair of hitters. We analyzed this question by tracking hitter pairs' season-long and pair-specific production at the PA and pitch levels to see its impact on each hitter's run expectancy in 2025. In this study, we were unable to replicate the 2025 findings using 2026 data, though we did gain some interesting insights. Firstly, we discovered that the focal run expectancy mechanically absorbed much of the PA-level production statistics, as shown with our default conditional OBP and focal SLG relationship. Once we pivoted to a proxy, focal Barrel/PA, we came across an interesting pattern. Two of our regressions: (3) Conditional Sweet Spot\% x Focal Attack Angle and (4) Line-Drive\% x Swing Path Tilt, had interaction coefficients that flipped signs, speaking to a potential issue beyond just our initial dependent variable. As we discovered with our data, a small number of frequently paired pairings represented a large contingent of our data, so expanding the dataset to include seasons going back as far as 2015 (the initial season of Statcast's batted ball metrics) could help gain more frequent pairs that can make our dataset more diverse for analysis.

Our investigation into Barrel/PA as a proxy suggests there may be a better way to understand how hitters are affected by those in front of or behind them. How do we enhance a player's power? Their on-base ability? What if there is a particular conditional player skill that keeps a focal high-chase hitter facing more pitches in the zone? Rather than focusing only on run output, hitter pairings could be established to enhance the player's output by highlighting their best skill.

Another potential avenue for this research is to assess how rigid pitchers are from at-bat to at-bat. An underlying theme of our framework was how pitchers might adjust from one hitter to the next in determining how that affects each hitter's run expectancy, particularly with our regression using Nitro Dist. For example, say a pitcher is facing a weak hitter and challenges him with pitches in the middle of the plate to get him out. How long does the pitcher need then to hone in his command of the edges of the zone? The number of batters it takes to adjust could be the difference between that pitcher staying in the game or not.

\newpage
\bibliographystyle{plainnat}
\bibliography{sample}

\appendix

\newpage
\section{Supplementary Tables}\label{app:tables}

\tbl{expected_signs_focal}
\tbl{expected_signs_cond}
\tbl{interaction_focal_2025}
\tbl{interaction_cond_2025}
\tbl{interaction_focal_2026}
\tbl{interaction_cond_2026}
\tbl{interaction_barrel_2025}
\tbl{interaction_barrel_2026}

\end{document}
